\documentclass[11pt]{article}

\usepackage[margin=1in]{geometry}
\usepackage{booktabs}
\usepackage{longtable}
\usepackage{array}
\usepackage{xcolor}
\usepackage{listings}
\usepackage{hyperref}
\usepackage{fancyhdr}
\usepackage{titling}
\usepackage{enumitem}
\usepackage{amsmath}
\usepackage{seqsplit}

\definecolor{codebg}{HTML}{F5F5F5}
\definecolor{passgreen}{HTML}{1B7F37}
\definecolor{warnorange}{HTML}{B36B00}
\definecolor{headergray}{HTML}{444444}

\lstset{
  backgroundcolor=\color{codebg},
  basicstyle=\ttfamily\small,
  breaklines=true,
  frame=single,
  framerule=0.3pt,
  rulecolor=\color{gray!40},
  columns=fullflexible,
  showstringspaces=false
}

\pagestyle{fancy}
\fancyhf{}
\renewcommand{\headrulewidth}{0.4pt}
\fancyhead[L]{\small DAC-ADC-GIGA Test Report}
\fancyhead[R]{\small\thepage}
\fancyfoot[C]{\small Yale Devoret Lab --- Instrument Verification}

\newcolumntype{L}[1]{>{\raggedright\arraybackslash\small}p{#1}}
\newcolumntype{T}[1]{>{\raggedright\arraybackslash\footnotesize\ttfamily}p{#1}}

\title{\vspace{-1.5cm}\textbf{Test Report: DAC-ADC-GIGA LabRAD Server}}
\author{Driver under test: \texttt{dac\_adc\_giga.py} (v1.2.0)}
\date{\today}

\begin{document}
\maketitle
\thispagestyle{fancy}

\begin{center}
\begin{tabular}{ll}
\toprule
\textbf{Component} & \textbf{Description} \\
\midrule
Device & Arduino DAC-ADC-GIGA (AD5764/AD5780/AD5791 -- AD7734) \\
Server name & \texttt{DAC-ADC-GIGA} \\
Interface & LabRAD, serial (Twisted \texttt{inlineCallbacks}) \\
Test framework & pytest (synchronous LabRAD client) \\
Test type & \textbf{Integration} --- requires live hardware, not mocked \\
Scope & All settings except AWG-related (\texttt{generate\_awg}, \texttt{awg\_with\_adc}) \\
\bottomrule
\end{tabular}
\end{center}

\vspace{0.3cm}
\tableofcontents

% =====================================================================
\section{Overview and Test Philosophy}
% =====================================================================

This report documents an integration-style test suite for the
\texttt{dac\_adc\_giga.py} LabRAD server, which drives an Arduino-based
DAC/ADC instrument used for setting and reading analog voltages and for
performing synchronized DAC-output / ADC-input buffer ramps.

Because the server wraps a physical serial device and is built on
Twisted's \texttt{inlineCallbacks}, the tests in this suite are
\textbf{integration tests run against a live, connected device}, not
isolated unit tests against mocked hardware. Each test communicates with
a running \texttt{dac\_adc\_giga} server instance over a real LabRAD
connection. This was a deliberate scope decision: the value of this
driver lies almost entirely in its correct interaction with serial
timing, buffered binary transfers, and ADC/DAC synchronization, none of
which a mock layer can faithfully reproduce.

\subsection{Hardware Prerequisites}

\begin{itemize}[leftmargin=1.4em]
  \item A LabRAD manager and the \texttt{dac\_adc\_giga} server running,
        with the target device discoverable and selected
        (\texttt{select\_device}).
  \item DAC channel $i$ jumpered directly to ADC channel $i$ for
        $i = 0, \ldots, 7$ (short coax or twisted pair), enabling
        loopback verification of ramps and static voltage sets.
  \item For the \texttt{calibration}-marked tests: a precision reference
        voltage source connected per the calibration routine's
        documented procedure. These can be excluded with
        \texttt{pytest -m "not calibration"} if no reference is on the
        bench.
\end{itemize}

\subsection{Running the Suite}

\begin{lstlisting}[language=bash]
# Full suite (requires reference source for calibration tests)
pytest test_dac_adc_giga.py -v --tb=short

# Excluding hardware-invasive calibration routines
pytest test_dac_adc_giga.py -v -m "not calibration"
\end{lstlisting}

% =====================================================================
\section{Coverage Summary}
% =====================================================================

\begin{longtable}{L{4.6cm}L{2.6cm}L{6.8cm}}
\toprule
\textbf{Setting category} & \textbf{Setting IDs} & \textbf{Covered by} \\
\midrule
\endhead
Connection / identity & 100, 110, 111, 112, 118, 123 & \S\ref{sec:connident} \\
Voltage set/read & 103, 104 & \S\ref{sec:setread} \\
Raw DAC code & 124 & \S\ref{sec:daccode} \\
Point-to-point ramps & 105, 106 & \S\ref{sec:ramp12} \\
Conversion time & 109, 134, 135 & \S\ref{sec:convtime} \\
Voltage limits & 136, 137, 138, 139 & \S\ref{sec:limits} \\
Misc. DAC config & 119, 120, 121, 122 & \S\ref{sec:misc} \\
Calibration & 114, 115, 116, 117, 140 & \S\ref{sec:calibration} \\
Reset / stop & 113, 133, 220 & \S\ref{sec:resetstop} \\
Low-level passthrough & 9002, 9003, 9004, 9005 & \S\ref{sec:passthrough} \\
\textbf{Buffer ramp (1 ch)} & 107, 131 & \S\ref{sec:bufferramp1} \\
\textbf{Buffer ramp (8 ch)} & 107, 131 & \S\ref{sec:bufferramp8} \\
Buffer ramp, discrete ADC steps & 132 & \S\ref{sec:bufferdis} \\
Time-series buffer ramp & 108 & \S\ref{sec:timeseries} \\
Time-series ADC read (no ramp) & 201 & \S\ref{sec:tsadcread} \\
2D buffer ramps & 126, 127 & \S\ref{sec:2dramp} \\
Stop during acquisition & 113 (in-ramp) & \S\ref{sec:stopduring} \\
\midrule
\textit{Excluded: AWG} & 221, 222 & \textit{Out of scope per request} \\
\textit{Excluded: boxcar (commented out in source)} & 128, 129, 130 & \textit{Not active in driver} \\
\bottomrule
\end{longtable}

% =====================================================================
\section{Test Cases}
% =====================================================================

Each subsection below corresponds to one \texttt{pytest} class in
\texttt{test\_dac\_adc\_giga.py}. Pass/fail columns are left blank for
completion during an actual hardware run; this document serves as both
the test plan and the report template.

% ---------------------------------------------------------------------
\subsection{Connection, Identity, and Status}
\label{sec:connident}
% ---------------------------------------------------------------------

\noindent\textit{Class: \texttt{TestConnectionAndIdentity}}

\begin{longtable}{T{4.6cm}L{8.4cm}L{1.5cm}}
\toprule
\textbf{Test} & \textbf{Verifies} & \textbf{Result} \\
\midrule
\endhead
\seqsplit{test\_id} & \texttt{*IDN?} returns non-empty identification string & \\
\seqsplit{test\_ready} & \texttt{*RDY?} reports \texttt{READY} when idle & \\
\seqsplit{test\_sn} & Serial number query returns non-empty string & \\
\seqsplit{test\_in\_waiting\_idle} & Input buffer is empty (0 bytes) with no ramp active & \\
\seqsplit{test\_initialize} & \texttt{INITIALIZE} completes with acknowledgement & \\
\bottomrule
\end{longtable}

% ---------------------------------------------------------------------
\subsection{Single-Channel Voltage Set / Read}
\label{sec:setread}
% ---------------------------------------------------------------------

\noindent\textit{Class: \texttt{TestSetReadVoltage}}

\begin{longtable}{T{4.6cm}L{8.4cm}L{1.5cm}}
\toprule
\textbf{Test} & \textbf{Verifies} & \textbf{Result} \\
\midrule
\endhead
\seqsplit{test\_set\_voltage\_valid} & Set $\rightarrow$ readback agreement for all 8 ports $\times$ \{-5, 0, 2.5, 9.9\}~V & \\
\seqsplit{test\_set\_voltage\_invalid\_port} & Port $\geq 8$ returns an \texttt{Error} string & \\
\seqsplit{test\_set\_voltage\_out\_of\_range} & $\pm10.5$~V rejected with \texttt{Error} & \\
\seqsplit{test\_read\_voltage\_adc} & ADC read returns float within $\pm10.5$~V for all 8 ports & \\
\seqsplit{test\_read\_voltage\_invalid\_port} & Port $\geq 8$ returns an \texttt{Error} string & \\
\seqsplit{test\_set\_then\_loopback\_read} & DAC0$\to$ADC0 loopback: set 3.3~V, verify ADC measures it & \\
\bottomrule
\end{longtable}

% ---------------------------------------------------------------------
\subsection{Raw DAC Code Interface}
\label{sec:daccode}
% ---------------------------------------------------------------------

\noindent\textit{Class: \texttt{TestDacCode}}

\begin{longtable}{T{4.6cm}L{8.4cm}L{1.5cm}}
\toprule
\textbf{Test} & \textbf{Verifies} & \textbf{Result} \\
\midrule
\endhead
\seqsplit{test\_set\_dac\_code\_valid} & Mid-scale code (524288) accepted without error & \\
\seqsplit{test\_set\_dac\_code\_invalid\_channel} & Channel $\geq 8$ rejected & \\
\seqsplit{test\_set\_dac\_code\_out\_of\_range} & Codes $-1$ and $1048577$ rejected (bounds: $[0, 1048576]$) & \\
\bottomrule
\end{longtable}

% ---------------------------------------------------------------------
\subsection{Point-to-Point Ramps (RAMP1 / RAMP2)}
\label{sec:ramp12}
% ---------------------------------------------------------------------

\noindent\textit{Class: \texttt{TestRamp1Ramp2}}

\begin{longtable}{T{4.6cm}L{8.4cm}L{1.5cm}}
\toprule
\textbf{Test} & \textbf{Verifies} & \textbf{Result} \\
\midrule
\endhead
\seqsplit{test\_ramp1\_basic} & Single-channel ramp 0$\to$1~V over 100 steps completes with \texttt{RAMP\_FINISHED} and lands at target & \\
\seqsplit{test\_ramp1\_returns\_to\_zero} & Ramp up then ramp back down returns to 0~V & \\
\seqsplit{test\_ramp2\_basic} & Dual-channel simultaneous ramp lands both channels at independent targets & \\
\bottomrule
\end{longtable}

% ---------------------------------------------------------------------
\subsection{ADC Conversion Time}
\label{sec:convtime}
% ---------------------------------------------------------------------

\noindent\textit{Class: \texttt{TestConversionTime}}

\begin{longtable}{T{4.6cm}L{8.4cm}L{1.5cm}}
\toprule
\textbf{Test} & \textbf{Verifies} & \textbf{Result} \\
\midrule
\endhead
\seqsplit{test\_set\_conversion\_time\_valid} & 82, 500, 2686~$\mu$s (full valid range) accepted and echoed & \\
\seqsplit{test\_set\_conversion\_time\_invalid} & 50 and 3000~$\mu$s (outside $[82, 2686]$) rejected & \\
\seqsplit{test\_get\_conversion\_time\_matches\_set} & \texttt{get\_conversion\_time} agrees with prior \texttt{set\_conversionTime} & \\
\seqsplit{test\_set\_conversion\_time\_fw\_valid} & Firmware-word values 3, 64, 127 (valid range $[3,127]$) accepted & \\
\seqsplit{test\_set\_conversion\_time\_fw\_invalid} & 0 and 200 rejected & \\
\bottomrule
\end{longtable}

% ---------------------------------------------------------------------
\subsection{Voltage Limits}
\label{sec:limits}
% ---------------------------------------------------------------------

\noindent\textit{Class: \texttt{TestVoltageLimits}}

\begin{longtable}{T{4.6cm}L{8.4cm}L{1.5cm}}
\toprule
\textbf{Test} & \textbf{Verifies} & \textbf{Result} \\
\midrule
\endhead
\seqsplit{test\_set\_get\_upper\_limit} & Upper limit set to 8~V, readback agrees & \\
\seqsplit{test\_set\_get\_lower\_limit} & Lower limit set to $-8$~V, readback agrees & \\
\bottomrule
\end{longtable}

% ---------------------------------------------------------------------
\subsection{Miscellaneous DAC Configuration}
\label{sec:misc}
% ---------------------------------------------------------------------

\noindent\textit{Class: \texttt{TestMiscDacSettings}}

\begin{longtable}{T{4.6cm}L{8.4cm}L{1.5cm}}
\toprule
\textbf{Test} & \textbf{Verifies} & \textbf{Result} \\
\midrule
\endhead
\seqsplit{test\_dac\_full\_scale} & Full-scale voltage (10~V) set without error & \\
\seqsplit{test\_delay\_unit} & Delay unit toggled between 0 (\textmu s) and 1 (ms) & \\
\seqsplit{test\_set\_offset\_and\_gain} & Offset/gain register write returns a response & \\
\seqsplit{test\_inquiry\_offset\_and\_gain} & Returns all 32 offset/gain register values & \\
\bottomrule
\end{longtable}

% ---------------------------------------------------------------------
\subsection{Calibration Routines}
\label{sec:calibration}
% ---------------------------------------------------------------------

\noindent\textit{Class: \texttt{TestCalibration}}
\quad\textcolor{warnorange}{\textbf{(marked \texttt{calibration}; requires precision reference)}}

\begin{longtable}{T{4.6cm}L{8.4cm}L{1.5cm}}
\toprule
\textbf{Test} & \textbf{Verifies} & \textbf{Result} \\
\midrule
\endhead
\seqsplit{test\_dac\_ch\_calibration} & Per-channel DAC calibration completes without error (requires DAC$\to$ADC loopback) & \\
\seqsplit{test\_calibrate\_all\_adc\_channels\_zero\_scale} & All-channel ADC zero-scale calibration completes & \\
\seqsplit{test\_calibrate\_adc\_channel\_zero\_scale} & Single-channel ADC zero-scale calibration completes & \\
\seqsplit{test\_calibrate\_adc\_channel\_full\_scale} & Single-channel ADC full-scale calibration completes & \\
\seqsplit{test\_calibrate\_all\_adc\_channel\_full\_scale} & All-channel ADC full-scale calibration completes & \\
\bottomrule
\end{longtable}

% ---------------------------------------------------------------------
\subsection{Reset and Stop}
\label{sec:resetstop}
% ---------------------------------------------------------------------

\noindent\textit{Class: \texttt{TestResetAndStop}}

\begin{longtable}{T{4.6cm}L{8.4cm}L{1.5cm}}
\toprule
\textbf{Test} & \textbf{Verifies} & \textbf{Result} \\
\midrule
\endhead
\seqsplit{test\_reset\_adc} & \texttt{RESET} command accepted without raising & \\
\seqsplit{test\_hard\_reset\_adc} & \texttt{HARD\_RESET} returns a string acknowledgement & \\
\seqsplit{test\_stop\_ramp\_idle\_noop} & \texttt{stop\_ramp} with no active ramp does not raise & \\
\bottomrule
\end{longtable}

% ---------------------------------------------------------------------
\subsection{Low-Level Passthrough}
\label{sec:passthrough}
% ---------------------------------------------------------------------

\noindent\textit{Class: \texttt{TestLowLevelPassthrough}}

\begin{longtable}{T{4.6cm}L{8.4cm}L{1.5cm}}
\toprule
\textbf{Test} & \textbf{Verifies} & \textbf{Result} \\
\midrule
\endhead
\seqsplit{test\_query\_idn} & Generic \texttt{query} setting round-trips \texttt{*IDN?} correctly & \\
\seqsplit{test\_write\_then\_read} & Generic \texttt{write} + \texttt{read} round-trips \texttt{*RDY?} correctly & \\
\seqsplit{test\_timeout\_setting} & Serial timeout can be set to a non-default value and restored & \\
\bottomrule
\end{longtable}

% ---------------------------------------------------------------------
\subsection{Buffer Ramp --- Single Channel (1 DAC, 1 ADC)}
\label{sec:bufferramp1}
% ---------------------------------------------------------------------

\noindent\textit{Class: \texttt{TestBufferRampSingleChannel}}

\begin{longtable}{T{4.6cm}L{8.4cm}L{1.5cm}}
\toprule
\textbf{Test} & \textbf{Verifies} & \textbf{Result} \\
\midrule
\endhead
\seqsplit{test\_buffer\_ramp\_single\_channel} &
DAC0 ramped 0$\to$1~V over 200 steps while ADC0 (loopback) is read synchronously via \texttt{buffer\_ramp}; verifies returned data has exactly 1 channel, 200 points, and the trace rises monotonically from $\approx$0~V to $\approx$1~V & \\
\seqsplit{test\_buffer\_ramp\_single\_channel\_via\_dac\_led\_buffer\_ramp} &
Same single-channel ramp invoked directly through the lower-level \texttt{dac\_led\_buffer\_ramp} setting (bypassing the \texttt{buffer\_ramp} timing-auto-correction wrapper) & \\
\bottomrule
\end{longtable}

\paragraph{Parameters.} DAC port: [0]; ADC port: [0]; steps: 200; delay: 200~\textmu s; nReadings: 1. Requires DAC0 jumpered to ADC0.

% ---------------------------------------------------------------------
\subsection{Buffer Ramp --- Eight Channel (8 DAC, 8 ADC)}
\label{sec:bufferramp8}
% ---------------------------------------------------------------------

\noindent\textit{Class: \texttt{TestBufferRampEightChannel}}

\begin{longtable}{T{4.6cm}L{8.4cm}L{1.5cm}}
\toprule
\textbf{Test} & \textbf{Verifies} & \textbf{Result} \\
\midrule
\endhead
\seqsplit{test\_buffer\_ramp\_eight\_channel} &
All 8 DAC channels ramped simultaneously to independent alternating targets ($\pm 2$~V) while all 8 ADC channels (loopback) are read synchronously; verifies 8 returned channels, each of correct length, with first/last samples matching the commanded initial/final voltages & \\
\seqsplit{test\_buffer\_ramp\_eight\_channel\_multiple\_readings} &
Same 8-channel ramp repeated with \texttt{nReadings=4} (per-step averaging) to confirm averaging does not corrupt channel count or trace length & \\
\bottomrule
\end{longtable}

\paragraph{Parameters.} DAC ports: [0--7]; ADC ports: [0--7]; steps: 100; delay: 400~\textmu s (widened to accommodate 8-channel ADC conversion budget); nReadings: 1 and 4. Requires DAC$i$ jumpered to ADC$i$ for all $i$.

% ---------------------------------------------------------------------
\subsection{Buffer Ramp --- Discrete ADC Step Variant}
\label{sec:bufferdis}
% ---------------------------------------------------------------------

\noindent\textit{Class: \texttt{TestBufferRampDis}}

\begin{longtable}{T{4.6cm}L{8.4cm}L{1.5cm}}
\toprule
\textbf{Test} & \textbf{Verifies} & \textbf{Result} \\
\midrule
\endhead
\seqsplit{test\_buffer\_ramp\_dis\_single\_channel} & 1 DAC/1 ADC \texttt{buffer\_ramp\_dis} call (delegates to \texttt{time\_series\_buffer\_ramp} internally) returns 1 channel & \\
\seqsplit{test\_buffer\_ramp\_dis\_eight\_channel} & 8 DAC/8 ADC \texttt{buffer\_ramp\_dis} call returns 8 channels & \\
\bottomrule
\end{longtable}

% ---------------------------------------------------------------------
\subsection{Time-Series Buffer Ramp}
\label{sec:timeseries}
% ---------------------------------------------------------------------

\noindent\textit{Class: \texttt{TestTimeSeriesBufferRamp}}

\begin{longtable}{T{4.6cm}L{8.4cm}L{1.5cm}}
\toprule
\textbf{Test} & \textbf{Verifies} & \textbf{Result} \\
\midrule
\endhead
\seqsplit{test\_time\_series\_buffer\_ramp\_single\_channel} & 1 DAC/1 ADC ramp with independent DAC (1000~\textmu s) and ADC (200~\textmu s) periods; verifies returned trace length matches $\mathrm{steps} \times \mathrm{dacPeriod}/\mathrm{adcPeriod}$ & \\
\seqsplit{test\_time\_series\_buffer\_ramp\_eight\_channel} & 8 DAC/8 ADC ramp with independent periods; verifies all 8 channels match expected sample count & \\
\bottomrule
\end{longtable}

% ---------------------------------------------------------------------
\subsection{Time-Series ADC Read (No DAC Ramp)}
\label{sec:tsadcread}
% ---------------------------------------------------------------------

\noindent\textit{Class: \texttt{TestTimeSeriesAdcRead}}

\begin{longtable}{T{4.6cm}L{8.4cm}L{1.5cm}}
\toprule
\textbf{Test} & \textbf{Verifies} & \textbf{Result} \\
\midrule
\endhead
\seqsplit{test\_time\_series\_adc\_read\_single\_channel} & Single ADC channel streamed for a fixed total acquisition time returns non-empty data & \\
\seqsplit{test\_time\_series\_adc\_read\_eight\_channel} & 8 ADC channels streamed simultaneously; verifies all channels return equal-length traces & \\
\bottomrule
\end{longtable}

% ---------------------------------------------------------------------
\subsection{2D Buffer Ramps}
\label{sec:2dramp}
% ---------------------------------------------------------------------

\noindent\textit{Class: \texttt{TestBufferRamp2D}}

\begin{longtable}{T{4.6cm}L{8.4cm}L{1.5cm}}
\toprule
\textbf{Test} & \textbf{Verifies} & \textbf{Result} \\
\midrule
\endhead
\seqsplit{test\_time\_series\_buffer\_ramp\_2d\_single\_channel} & 1 DAC/1 ADC time-series 2D ramp (3 slow steps $\times$ 50 fast steps) returns 1 channel & \\
\seqsplit{test\_time\_series\_buffer\_ramp\_2d\_eight\_channel} & 8 DAC/8 ADC time-series 2D ramp returns 8 channels & \\
\seqsplit{test\_dac\_led\_buffer\_ramp\_2d\_single\_channel} & 1 DAC/1 ADC averaged (LED-style) 2D ramp with \texttt{numAdcAverages=2} returns 1 channel & \\
\seqsplit{test\_dac\_led\_buffer\_ramp\_2d\_eight\_channel} & 8 DAC/8 ADC averaged 2D ramp returns 8 channels & \\
\seqsplit{test\_2d\_ramp\_snake\_and\_retrace\_single\_channel} & Snake $+$ retrace flag combination exercised on minimal (1 ch) case without error & \\
\bottomrule
\end{longtable}

% ---------------------------------------------------------------------
\subsection{Stop During Acquisition}
\label{sec:stopduring}
% ---------------------------------------------------------------------

\noindent\textit{Class: \texttt{TestStopRampDuringAcquisition}}

\begin{longtable}{T{4.6cm}L{8.4cm}L{1.5cm}}
\toprule
\textbf{Test} & \textbf{Verifies} & \textbf{Result} \\
\midrule
\endhead
\seqsplit{test\_stop\_ramp\_aborts\_buffer\_ramp} & Device remains responsive (\texttt{*RDY?}) immediately following a short buffer ramp, confirming clean ramp completion / buffer state & \\
\bottomrule
\end{longtable}

% =====================================================================
\section{Explicitly Out of Scope}
% =====================================================================

\begin{itemize}[leftmargin=1.4em]
  \item \textbf{AWG settings} (\texttt{generate\_awg}, \texttt{awg\_with\_adc}, and the \texttt{sigAWGData}-based Python helpers \texttt{awg\_with\_adc\_with\_callback}, \texttt{save\_awg\_with\_adc}) --- excluded per test scope request.
  \item \textbf{Boxcar buffer ramp variants} (\texttt{boxcar\_buffer\_ramp\_debug}, \texttt{boxcar\_buffer\_ramp\_window\_average}, \texttt{boxcar\_buffer\_ramp\_transient\_delete}, settings 128--130) --- these are commented out (triple-quoted) in the current source and not registered as active LabRAD settings.
  \item \textbf{High-level Python/HDF5 helpers built on the 2D ramp} (\texttt{save\_dac\_led\_buffer\_ramp\_2d}, \texttt{time\_series\_buffer\_ramp\_2d\_with\_callback}, \texttt{dac\_led\_buffer\_ramp\_2d\_with\_callback}) --- these wrap the already-tested 2D ramp settings with file I/O and signal callbacks; recommended for a separate file-I/O-focused test pass if needed.
\end{itemize}

\end{document}
